

Tropical forests are the world's key ecosystems providing habitat for a wide diversity of species, sustaining local livelihoods and regulating global ecological cycles. Over half of tropical forests have been transformed into other land uses, mainly agriculture. Shifting agriculture, a common agricultural practice in the tropics, often results in land abandonment which facilitates natural regeneration of forests. Given the growing pressure on tropical forests, natural forest regeneration represents a promising solution for ecological restoration. In this study, we analyzed the impact of surrounding forest landscape on the seed rain of woody plants in naturally regenerating secondary tropical forests. We collected seeds using seed traps in early successional plots across wet and dry tropical forests in southern Mexico. All seeds were identified to species level and classified by dispersal mode (biotic/abiotic) and functional guild (pioneer, shade-tolerant, generalist). To characterize the surrounding forest landscape, we used Landsat satellite time series (1992–2022) and the AVOCADO algorithm to map forest cover, connectivity, and successional stage within a 1 km buffer around each plot. We then applied correlation analyses and regression models to test how these forest landscape variables influenced seed rain richness, dispersal mode, and guild composition across forest types. Seed richness was significantly higher in wet forests compared to dry forests. In both forest types, higher late successional forest cover was associated with greater seed richness. Dispersal mode differed strongly between forest types: wet forests were dominated by biotically dispersed seeds, while dry forests had a greater proportion of abiotically dispersed seeds. Guild composition showed weaker responses, although proportion of generalist species decreased with increasing forest cover in both forest types. Our results suggest that older successional forests in the surrounding landscape are essential for enhancing seed richness. Furthermore, well-connected forests may be more important in wet tropical forests dominated by biotically dispersed seeds, whereas more open landscapes may play a greater role in dry tropical forests dominated by abiotically dispersed seeds. Therefore, considering the successional stage of the surrounding forest landscape and adopting region-specific conservation plans are pivotal when aiming to naturally regenerate degraded landscapes.

\textbf{Keywords:} seed rain, forest landscape, AVOCADO, remote sensing, restoration